\section{Legal Landscape and State Court Pathways}
\label{sec:legal}

\subsection{Post-\textit{Rucho} Federal Foreclosure}

\textit{Rucho v.\ Common Cause} (2019) forecloses federal judicial review of
partisan gerrymandering claims under the Constitution. The holding rests on
two grounds: (1) no judicially manageable standard exists to distinguish
permissible partisan line-drawing from unconstitutional excess, and (2) partisan
gerrymandering claims present political questions that federal courts lack the
institutional competence to resolve. The majority's survey of proposed standards---
equal protection, First Amendment association, Elections Clause---found each
either too broad (invalidating all partisan redistricting) or too indeterminate
(requiring courts to set the acceptable level of partisanship, a choice the
majority characterized as ``not law'' but politics).\footnote{\textit{Rucho},
588 U.S.\ at 720--21.}

This forecloses not just claims under the Constitution's partisan gerrymandering
doctrine but also the federal common law of redistricting as a check on partisan
maps. Federal courts can still review congressional redistricting for
\textit{Wesberry} population equality violations, \textit{Shaw v.\ Reno} racial
gerrymandering, and Voting Rights Act compliance---but not partisan composition.

\subsection{State Constitutional Pathways}

Three active state constitutional pathways have produced successful partisan
gerrymandering claims since \textit{Rucho}:

\textbf{Pennsylvania's Free and Equal Elections Clause} (Art.\ I, \S 5). The
Pennsylvania Supreme Court, in \textit{League of Women Voters v.\ Commonwealth}
(2018)---decided before \textit{Rucho}, but under a state constitutional theory
immune to federal foreclosure---held that Pennsylvania's congressional map
violated the free and equal elections clause. The court drew a remedial map
itself, applying compactness, contiguity, and population equality criteria without
partisan data. The EG and MM of the court-drawn remedial map were substantially
lower than the invalidated enacted map, demonstrating that state constitutional
standards can be operationalized through metric evidence.

\textbf{North Carolina's Free Elections Clause} (Art.\ I, \S 10). In
\textit{Harper v.\ Hall} (N.C.\ 2022), the North Carolina Supreme Court
held that the legislature ``may not draw districts for the purpose of
systematically advantaging or disadvantaging a particular political party''
and that the 2022 congressional map violated this standard. The court relied
on multiple partisan fairness metrics---including EG and MM---as evidence
that the legislature had subordinated all legitimate redistricting criteria
to partisan advantage. The court did not adopt a specific numerical threshold
for any metric; instead it treated the convergence of multiple metrics
pointing in the same direction as evidence of constitutional violation.

\textbf{Ohio's Anti-Gerrymandering Provision} (Ohio Const.\ Art.\ XI, \S 6).
The Ohio Constitution's 2015 amendment explicitly prohibits drawing district
boundaries ``for the purpose of favoring or disfavoring a political party.''
In \textit{League of Women Voters of Ohio v.\ Ohio Redistricting Commission}
(2022), the Ohio Supreme Court applied this standard to invalidate the Ohio
congressional map seven times over three years as the legislature repeatedly
submitted plans that failed the constitutional test. The court's analysis
included seat-share projections and vote-share analysis that map onto the
partisan metrics in this paper.

\subsection{The Metric Selection Problem in State Court Practice}

Expert witnesses in state redistricting litigation face a metric selection
problem: which of the six metrics provides the strongest, most defensible
evidence under the applicable state constitutional standard?

The answer depends on the legal theory. For NC's ``free elections'' clause,
which has been interpreted to prohibit systematic partisan advantage rather
than requiring proportionality, symmetry-based metrics (Partisan Bias, EG)
and the seats-votes curve are most directly relevant---they measure whether
the map treats the two parties symmetrically when their vote totals are equal.
For PA's free and equal elections clause, which has been interpreted to impose
a softer proportionality floor, the Gallagher Index and MM are also relevant.
For Ohio's explicit anti-gerrymandering provision, the full six-metric package
demonstrates systematic partisan intent most completely.

The decision table in Section~\ref{sec:definition} (Table~\ref{tab:decision-table})
operationalizes this guidance. Expert witnesses should lead with the metric
most directly tied to the applicable legal standard, report all six for
convergence, and explain that the \textsc{bisect} reference distribution
establishes the neutral baseline from which the enacted map's deviation
is measured. Metric-shopping criticism---the objection that the expert selected
the metric most favorable to their client---is defeated by reporting all six
simultaneously and showing that all six converge.

\subsection{What Courts Have and Have Not Adopted}

Through 2026, the following metric adoptions and rejections can be documented:

\begin{itemize}
\item \textbf{EG}: Cited in \textit{Gill v.\ Whitford} (federal, not adopted
  as standard), \textit{Whitford v.\ Gill} (7th Cir., applied but not dispositive),
  \textit{Harper v.\ Hall} (NC, treated as one of several indicators).
\item \textbf{MM}: Cited in \textit{Harper v.\ Hall} (NC), academic literature
  widely. No court has adopted a specific MM threshold.
\item \textbf{Partisan Bias}: Part of the symmetry framework cited in NC and PA
  expert testimony. Courts have treated it as relevant evidence but not dispositive.
\item \textbf{Declination}: Introduced as an exhibit in \textit{Common Cause v.\
  Rucho} (M.D.N.C.\ 2018, pre-\textit{Rucho} SCOTUS); the court did not rely on it.
  No federal or state court has adopted declination as evidence as of 2026 (see L.4).
\item \textbf{Seats-votes curve}: Admitted as expert evidence in NC and OH
  proceedings. Courts have used the graphical exhibit without adopting a specific
  responsiveness threshold.
\item \textbf{Gallagher LSq}: Used in academic analysis of redistricting; not
  adopted as a standalone legal standard in any U.S.\ court as of 2026.
\end{itemize}
